\chapter{Один tap}\label{ch:one-tap}
Касание карты и~звук терминала отделяет немногим больше одной секунды.
За~это время одиннадцать участников по~очереди читают и~переписывают одно
платёжное сообщение, и~каждый из~них живёт по~своим правилам, описанным в~разных
главах книги.

Предыдущие главы разбирали EMV, бесконтактный протокол, ISO~8583, криптографию,
ОПКЦ и~эквайера по~отдельности. Каждый разбор закрыт на~себе и~останавливается
там, где начинается чужая тема: гл.~\ref{ch:emv} не~доходит до~ОПКЦ,
гл.~\ref{ch:mir} не~разбирает GPO. Эта глава рисует шов между ними~--- как одна
APDU превращается в~ISO~8583 и~как ответ возвращается обратно в~карту; за~деталями
каждой фазы~--- ссылки на~исходные главы.

\section{Условия примера}\label{sec:one-tap-setup}

Конкретная транзакция: бесконтактная карта \enquote{Мир Дебетовая}, эмитент~---
АО~\enquote{Т-Банк}, эквайер~--- ПАО~Сбербанк, терминал~---
бесконтактный ридер на~кассе федеральной розничной сети, сумма
1\,420{,}50~руб. Карта приложена в~11:43:07 по~московскому времени, чек
распечатан в~11:43:09. Между этими двумя секундами умещаются одиннадцать фаз.
Их типовой бюджет~--- около 1200~мс, из~которых половина уходит на~обмен данными
карты с~терминалом, а~вторая половина~--- на~ОПКЦ и~эмитента.

Транзакция онлайновая: бесконтактный лимит CVM превышен, терминал обязан
запросить авторизацию по~правилам ПС~Мир. Это означает, что ARQC
обязательно уходит к~эмитенту и~beep терминала откладывается до~получения
ответа. Низкие суммы с~офлайн-подтверждением (\textit{tap-and-go})
дают другую последовательность~--- беззвучное снятие в~поле без~ожидания
сети, разобранное в~разделе~\ref{sec:contactless-emv}.

Весь внутрироссийский карточный трафик идёт через ОПКЦ НСПК с~1~апреля
2015~года~\cite{nspk-about-company} и~через ОПКЦ ежедневно проходит до~20~млн
операций~\cite{cnews-nspk-opkc}. 23~февраля 2024~года OFAC внёс
НСПК в~SDN-список, дополнительно ограничив зарубежный приём карт российских
эмитентов независимо от~бренда. Типовая платёжная транзакция в~России в~2026~--- это
карта \enquote{Мир}, российские эквайер и~эмитент, маршрут целиком внутри страны.

\begin{figure}[htbp]
  \centering
  \includefiguresvg[width=\linewidth]{assets/figures/ch13-one-tap/one-tap-budget}
  \caption{Расходование бюджета tap-to-beep: 1170~мс.}%
  \label{fig:one-tap-budget}
\end{figure}
\FloatBarrier

\section{Карта в~поле: фазы 1--5}\label{sec:one-tap-card-terminal}

Первые пять фаз~--- обмен бесконтактной карты и~терминала. Карта попадает
в~рабочую зону антенны, получает питание от~поля, исполняет EMV-диалог и~отдаёт
ARQC. Терминал в~это время инициализирует радиоинтерфейс, выбирает приложение,
читает запись карты и~получает криптограмму. По~независимым измерениям
бесконтактного EMV на~реальной терминальной шине этот обмен укладывается
в~500--600~мс, из~которых самая большая доля~--- одна
GENERATE~AC~\cite{lifecycleintegrity-emv-timing}.

\subsection*{Фаза~1. RF-активация, antiselect, T=CL handshake (0--25~мс)}

Терминал излучает поле на~13{,}56~МГц. Антенна карты ловит его, выпрямительный
диод заряжает накопительный конденсатор, чип запускается за~единицы миллисекунд.
Дальше идёт ISO~14443: WUPA (Wake-Up Type~A) от~терминала, ATQA (Answer To
Request, Type~A) от~карты, антиколлизия по~UID, SAK (Select Acknowledge),
RATS (Request for Answer To Select), ATS (Answer To Select). К~исходу этой
фазы у~карты и~терминала есть согласованный транспортный канал T=CL
по~ISO~14443-4~\cite{iso-14443-3,iso-14443-4}, с~максимальным временем ожидания
ответа FWT, по~умолчанию 4{,}8~мс при~FWI=4 (Frame Waiting Integer) и~до~4{,}9~с
при~FWI=14~\cite{iso-14443-4}.

ISO~14443 не~зависит от~карточной схемы: одна и~та же физика общая
для~Мир, Visa и~UnionPay. Подробно об~этом~--- раздел~\ref{sec:contactless-nfc}.

\subsection*{Фаза~2. PPSE и~выбор приложения (25--65~мс)}

Поверх T=CL терминал отправляет APDU \texttt{SELECT 2PAY.SYS.DDF01}~--- запрос
PPSE (Proximity Payment System Environment). Карта возвращает FCI, в~теге~A5
которого содержится список Application Templates (тег~61): для~каждого приложения~---
ADF Name (4F), Application Label (50) и~Application Priority Indicator
(87)~\cite{emvco-contactless-book-b}. На~карте \enquote{Мир} ADF Name начинается
с~RID НСПК \texttt{A0~00~00~06~58}, зарегистрированного по~ISO/IEC~7816-5
за~НСПК как application provider~\cite{iso-7816-5}.

Combination Selection пересекает AID карты со~списком ядер, которые поддерживает
терминал. В~ПС~Мир действует специальное правило: если на~карте
одновременно есть AID Мир и~AID другой системы (например, в~кобейджинговой
конструкции Мир--UnionPay), терминал в~России выбирает только AID~Мир и~передаёт
управление Mir~Kernel~\cite{nspk-mir-rules-v43,nexo-mir-fast-2022}.
Mir~Kernel~--- собственное бесконтактное ядро НСПК; его терминальный интерфейс
описан в~публикуемой НСПК спецификации \enquote{Mir Payment System Contactless
Terminal Kernel Specification}~\cite{nspk-mir-kernel}, низкоуровневые детали
криптограммных проверок остаются в~закрытой части документации для~участников.

Карта приоритетов разобрана в~разделе~\ref{sec:emv-app-selection};
кобейджинг Мир--UnionPay и~его исходы~--- в~разделе~\ref{sec:mir-opkc}.

\subsection*{Фаза~3. GPO, AIP, AFL (65--110~мс)}

Терминал отправляет \texttt{GET~PROCESSING~OPTIONS} с~PDOL, в~который запакованы
сумма, валюта, страна терминала, capabilities, TTQ и~UN. Карта возвращает AIP
(Application Interchange Profile~--- какие проверки приложение поддерживает)
и~AFL (Application File Locator~--- адреса записей, которые терминал должен
прочитать). Структура AIP\,/\,AFL~--- общая EMV-семантика, одинаковая по~Books
EMVCo и~по~Mir~Kernel.

PDOL, AIP и~AFL разобраны в~разделе~\ref{sec:emv-architecture}.

\subsection*{Фаза~4. READ~RECORD и~аутентификация данных (110--310~мс)}

По~AFL терминал читает несколько SFI/Record-записей: PAN (тег~5A) из~диапазона
\texttt{2200--2204} по~ISO/IEC~7812~\cite{iso-7812-transition}, expiry (5F24),
CDOL1, сертификат эмитента (90), открытый ключ эмитента (92) и~данные карты
для~аутентификации. На~картах \enquote{Мир} с~криптографией по~ГОСТ
(см.~раздел~\ref{sec:crypto-gost}) сертификатная иерархия и~алгоритмы подписи
строятся на~ГОСТ~34.10--2018 и~ГОСТ~34.11--2018; на~картах \enquote{Мир}
с~международной EMV-криптографией~--- на~RSA\,/\,SHA. Терминальное ядро
валидирует цепочку по~корневым ключам CA, зашитым через персонализацию
TMS~\cite{nspk-mir-kernel}.

Иерархия CA\,/\,IPK\,/\,ICC и~общая схема CDA\,/\,fDDA~---
разделы~\ref{sec:emv-data-auth} и~\ref{sec:crypto-key-hierarchy}.

\subsection*{Фаза~5. GENERATE~AC и~ARQC (310--595~мс)}

Терминал собирает CDOL1 и~отправляет \texttt{GENERATE~AC} с~P1=\texttt{0x80}~---
запрос ARQC (Authorisation Request Cryptogram). Карта внутри Secure Element
выводит сессионный ключ SK\textsubscript{AC} из~UDK по~схеме, согласованной с~НСПК
(для~ГОСТ-карт ключ выводится через Кузнечик\,/\,Магма; для~карт с~международной
криптографией~--- по~стандартной EMV-схеме на~3DES, см.~раздел~\ref{sec:crypto-key-hierarchy}),
считает MAC поверх данных транзакции, ATC и~криптограммных параметров и~возвращает
ARQC, ATC, IAD и~CID. По~измерениям бесконтактного EMV эта одна APDU занимает
порядка 280~мс на~пластиковой карте~\cite{lifecycleintegrity-emv-timing}~---
самая длинная фаза разговора карты с~терминалом, и~именно она задаёт нижний
предел \textit{tap-to-beep} даже при~идеальной сети. ГОСТ-карта в~теории может работать
медленнее из-за программной реализации Кузнечика на~отдельных чипах,
но~для~массовой эмиссии бюджет всё равно укладывается в~300~мс.

К~концу фазы~5 карта может покинуть поле: терминал собирает всё, что ему
нужно для~онлайн-авторизации, и~дальнейший разговор уже не~требует физического
присутствия карты. Бесконтактный Mir~Kernel почти никогда не~использует второй
\texttt{GENERATE~AC}: к~моменту ответа эмитента карта обычно уже в~кошельке.

Внутреннее устройство ARQC, набор полей CDOL1, разница ARQC\,/\,TC\,/\,AAC
и~место CID в~решении~--- раздел~\ref{sec:emv-cryptogram}.

\section{Сеть в~воздухе: фазы 6--10}\label{sec:one-tap-network}

После того как ARQC получен, начинается путешествие по~сети: терминал
упаковывает данные в~ISO~8583, передаёт эквайеру, эквайер маршрутизирует
в~ОПКЦ, ОПКЦ доставляет сообщение эмитенту, эмитент верифицирует криптограмму,
принимает решение, генерирует ARPC и~отправляет ответ обратно. Целиком внутри России, около 580~мс типового бюджета.

\subsection*{Фаза~6. Упаковка в~ISO~8583 0100 по~профилю НСПК (595--610~мс)}

Терминальное приложение собирает авторизационное сообщение в~формате
ISO~8583~\cite{iso-8583-2023}: MTI~0100 (Authorization Request)
с~bitmap'ом, указывающим присутствие полей DE~2 (PAN из~диапазона
\texttt{2200--2204}), DE~3 (Processing Code), DE~4 (Amount), DE~7, DE~11, DE~12,
DE~14, DE~22 (POS Entry Mode~--- бесконтактный), DE~25 (POS Condition Code),
DE~48 (Additional Data) и~DE~55 (Integrated Circuit Card System Related Data).
DE~55~--- это BER-TLV-капсула с~теми же EMV-тегами, что считала карта: ARQC
(9F26), ATC (9F36), IAD (9F10), AIP (82), AFL (94), CVR (Card Verification
Results) внутри IAD, и~так далее.
Полный набор полей и~их интерпретация определяются спецификацией НСПК и~доступны
только участникам ПС~Мир~\cite{nspk-about-company}; на~уровне абстракции
DE~55 переносит криптограмму и~контекст её формирования к~эмитенту без~потерь.

Сама упаковка занимает единицы миллисекунд: это работа CPU терминала, не~ввод-вывод.
Структура bitmap и~ISO~8583 в~целом~--- раздел~\ref{sec:iso8583-structure};
маршрут через ОПКЦ и~правила НСПК~--- раздел~\ref{sec:mir-opkc}.

\subsection*{Фаза~7. Терминал → FEP эквайера → шлюз к~ОПКЦ (610--680~мс)}

Терминал передаёт сообщение через тот канал, какой у~него настроен:
GPRS/4G через встроенный модем, проводной Ethernet через кассу, реже WiFi.
У~российских ТСП большинство стационарных терминалов работают через
выделенную сеть эквайера~--- VPN поверх 4G или MPLS-канал ТСП. Транспорт~--- TLS
с~клиентским сертификатом, который TMS прошивает в~терминал при~регистрации
(раздел~\ref{sec:ia-acquirer} разбирает TMS и~жизненный цикл прошивок).

На~входе у~эквайера сообщение принимает Front-End Processor (FEP)~--- сетевой
шлюз, который терминирует TLS, парсит ISO~8583, валидирует MAC терминала
(если используется session key поверх T=CL и~отдельный MAC хоста), и~по~BIN
определяет дальнейший маршрут. У~Сбера, как у~крупнейшего эквайера российского
рынка с~долей более 60\,\% по~количеству терминалов~\cite{comnews-sber-acquirer-2025},
FEP-инфраструктура построена на~собственных мощностях; у~небольших процессоров
FEP может быть арендован у~интегратора. BIN-таблица обновляется у~эквайера
из~централизованного перечня НСПК. По~BIN~\texttt{220070} (один из~диапазонов
Т-Банка) FEP понимает~--- карта \enquote{Мир}, маршрут на~шлюз эквайера к~ОПКЦ.

\emph{Шлюз эквайера к~ОПКЦ}~--- это та же роль, что у~Mastercard
Interface~Processor (MIP) в~международной сети, но~внутри российской
сети: терминирующая точка авторизационного протокола между процессором
и~национальным коммутатором. Технически это собственная разработка эквайера
или решение от~интегратора (например, \enquote{Инфосистемы Джет} работали с~НСПК
над~клиринговой системой Мир~\cite{tadviser-jet-mir-clearing}); внутреннего общего
\enquote{железного} стандарта на~аналог MIP в~НСПК нет.

Бюджет фазы~--- 50--80~мс по~российской магистральной сети, до~150~мс в~регионе
со~слабым каналом.

\subsection*{Фаза~8. ОПКЦ маршрутизация → шлюз эмитента (680--740~мс)}

ОПКЦ (операционный платёжный и~клиринговый центр НСПК)~--- единый центральный
коммутатор внутрироссийских карточных операций. В~отличие от~распределённой
сети Mastercard с~региональными Authorization Platform и~глобальной
маршрутизацией, ОПКЦ~--- один логический узел внутри России.
ОПКЦ идентифицирует эмитента по~BIN из~собственной таблицы маршрутизации,
формирует адрес шлюза эмитента и~доставляет сообщение
дальше~\cite{nspk-about-company,cnews-nspk-opkc}.

На~хорошей сети этот переход занимает 30--50~мс: одна короткая магистраль
без~глобальной маршрутизации между континентами. Сравнение с~Mastercard:
сигнал между ЦОДами Visa или Mastercard на~разных континентах не~уложится
меньше чем в~140~мс на~round trip; внутрироссийский Мир-маршрут физически
не~бывает таким длинным.

В~сценариях \textit{stand-in} ОПКЦ принимает решение об~авторизации вместо
недоступного эмитента по~параметрам, согласованным заранее. В~Правилах ПС~Мир
этот механизм называется \emph{Сервисом Резервной Авторизации}
(раздел~13)~\cite{nspk-mir-rules-v43}; международный аналог~--- Visa STIP
или Mastercard stand-in. См.~раздел~\ref{sec:mir-opkc}.

\subsection*{Фаза~9. Auth host эмитента: HSM verify ARQC, risk, ARPC (740--1020~мс)}

В~ЦОДе Т-Банка сообщение принимает Authorization Host~--- центральный процесс
эмитента, обслуживающий онлайн-авторизацию (раздел~\ref{sec:ia-issuer}).
Он распаковывает ISO~8583, достаёт из~DE~55 теги ARQC, ATC, IAD, AIP и~отдаёт
их HSM на~проверку.

HSM~--- у~российских эмитентов это либо сертифицированный
Thales~payShield с~прошивкой ГОСТ-модуля, либо отечественный HSM с~поддержкой
ГОСТ-2018 (Кузнечик, Магма, Стрибог)~--- получает команду на~верификацию ARQC.
По~серии команд HSM восстанавливает UDK из~IMK эмитента и~PAN, выводит
сессионный ключ SK\textsubscript{AC} по~схеме НСПК и~заново считает MAC поверх
предъявленных данных транзакции, сверяя результат с~ARQC карты. Один вызов HSM~---
единицы миллисекунд при~правильной конфигурации очередей.

Параллельно с~криптопроверкой работает риск-движок: лимиты по~карте, скоринг
по~MCC, поведенческие проверки, кэш статусов авторизаций, белые и~чёрные списки
устройств. Решение~--- approve, decline или soft-decline. В~бесконтактном
пути \textit{step-up} невозможен: карта не~в~поле, повторного \texttt{GENERATE~AC}
не~будет, поэтому soft-decline превращается в~обычный decline или approve.

Если решение положительное, HSM считает ARPC: MAC поверх ARQC и~ARC
(Authorisation Response Code) по~схеме, согласованной с~НСПК. ARPC возвращается
в~Authorization~Host. Типичный auth host решает за~150--300~мс по~отраслевым сводкам;
выделенные системы доводят это число до~единиц миллисекунд на~своей
стороне~\cite{qonto-card-auth}.

\subsection*{Фаза~10. ISO~8583 0110 и~обратный путь (1020--1100~мс)}

Authorization~Host формирует ответное сообщение MTI~0110: bitmap, тот же
DE~2\,/\,DE~3\,/\,DE~4 для~корреляции, DE~38 (Authorization ID Response Code),
DE~39 (Response Code: \texttt{00}~--- approved), DE~55 с~ARPC в~теге~91.
Сообщение уходит на~шлюз эмитента к~ОПКЦ, оттуда~--- в~ОПКЦ, оттуда~--- на~шлюз
эквайера, оттуда~--- в~FEP эквайера, оттуда~--- к~терминалу.

Обратный путь идёт по~тем же физическим узлам, что и~прямой, и~стоит примерно
столько же~--- около 80~мс. Коды отклика ISO~8583 и~правила интерпретации
DE~39~--- раздел~\ref{sec:iso8583-hex-dump}.

\section{ECR, чек, ОФД: фаза 11}\label{sec:one-tap-ecr}

После того как терминал получил MTI~0110 с~\texttt{Response~Code~00}, остаётся
последняя миля~--- сообщить о~результате кассе, напечатать чек и~отправить
фискальные данные оператору фискальных данных (ОФД).

\subsection*{Фаза~11. Терминал → ECR → касса → ОФД (1100--1170~мс)}

Между терминалом и~кассовой программой (ECR~--- Electronic Cash Register)
существует отдельный интерфейсный протокол. Терминал ничего не~знает о~корзине
покупателя и~артикулах; ECR ничего не~знает об~EMV и~ISO~8583. Между ними~---
один запрос-ответ по~локальной шине: \enquote{сделай платёж на~такую-то сумму}~---
\enquote{платёж сделан, RRN такой-то, код авторизации такой-то}.

Этот протокол существует в~двух мирах:

\emph{Стандартизированный nexo Retailer Protocol} (изначально продвигаемый
European Payments Council и~nexo standards, сейчас независимая организация),
построенный на~ISO~20022 и~передающий запрос\,/\,ответ как XML или
JSON-сообщения~\cite{nexo-retailer-protocol}. Цель nexo~--- развязать ТСП
от~конкретного эквайера: одно и~то же кассовое ПО работает с~любым терминалом,
любого поставщика, у~любого PSP. В~2022~году НСПК интегрировал ПС~Мир
в~nexo~FAST~--- ускоренный профиль nexo для~бесконтактных платежей~\cite{nexo-mir-fast-2022};
для~российского рынка это означает, что европейские интегрированные кассовые
решения могут принимать \enquote{Мир} без~отдельного драйвера эквайера.

\emph{Проприетарные протоколы PSP и~поставщиков терминалов}: каждый эквайер
исторически имел свой бинарный или текстовый формат связи между терминалом
и~кассой. Российские ТСП в~массе работают по~таким протоколам~--- свой
у~Сбера, свой у~ВТБ, свой у~Альфы, у~локальных поставщиков (\enquote{Эвотор},
АТОЛ)~--- с~интеграцией под~конкретное кассовое ПО (1С, R\,-Keeper, iiko,
Frontol).

После того как ECR получила подтверждение от~терминала, она формирует чек:
сумма, маскированный PAN, последние четыре цифры, дата, время, RRN, код
авторизации, имя ТСП, адрес. По~Федеральному закону от~22.05.2003~№~54-ФЗ
\enquote{О~применении контрольно-кассовой техники}~\cite{fz-54} к~этому
добавляются фискальные реквизиты~--- номер фискального накопителя (ФН), номер
фискального документа (ФД), фискальный признак документа (ФПД), QR-код для~проверки
в~приложении Федеральной налоговой службы (ФНС), и~всё это уходит на~ОФД через
интернет. ОФД пересылает данные в~ФНС.

Beep терминала и~печать чека выполняются параллельно: терминал издаёт звук
сразу по~получении \texttt{Response~Code~00}, кассовая программа печатает
несколько секунд. С~точки зрения покупателя \textit{tap-to-beep}~--- это около
1200~мс; \textit{tap-to-receipt}~--- ещё 2--4~секунды.

\section{Бюджет 1170~мс}\label{sec:one-tap-budget}

Из~1170~мс типового бесконтактного бюджета 595~мс приходится на~разговор карты
с~терминалом, 505~мс~--- на~сеть и~эмитента, 70~мс~--- на~ECR\,/\,печать чека
(см.~рис.~\ref{fig:one-tap-budget}).
Внутри карточно-терминальной половины 280~мс из~595~занимает одна
\texttt{GENERATE~AC}~--- криптограмма на~ограниченном CPU чипа. Внутри сетевой
половины 280~мс~--- Authorization~Host эмитента с~HSM-вызовами и~риск-движком.

Когда сумма ниже порога CVM и~эмитент допускает офлайн-одобрение
(\textit{tap-and-go}), фазы~7--10 пропадают целиком. Терминал принимает
решение по~правилам ядра, пока карта ещё в~поле, и~издаёт beep сразу после
фазы~5. Денежная сторона разрешается позже, через поздний клиринг (\textit{deferred
capture}): эквайер собирает одобренные офлайн транзакции пакетом, отправляет
их в~ОПКЦ, а~если на~счёте держателя средств не~хватит, разбирательство
случится постфактум. Tap-to-beep падает
до~600~мс~--- ровно вдвое. Именно эта схема обслуживает турникеты метро
и~автобусные валидаторы, где задержка в~полсекунды на~каждом пассажире
превращается в~очередь у~двери.

В~онлайн-сценарии главный соблазн оптимизации~--- разогнать криптографию:
ускорить HSM, поставить его в~один ЦОД с~auth host, сменить шифр. Из~280~мс
auth host эмитента около 270~мс занимает риск-движок~--- лимиты по~карте,
скоринг по~MCC, поведенческие проверки, кэши авторизаций и~списки устройств.
HSM на~verify~ARQC
и~генерации ARPC тратит единицы миллисекунд на~вызов; если HSM и~auth~host
находятся в~одном домене синхронной репликации (расстояние между ЦОДами
  до~50~км, ограничено скоростью света в~оптике и~оверхедом терминирующего
оборудования), один вызов укладывается в~2~мс. Колокация HSM с~auth~host
в~одном ЦОДе экономит меньше 2\,\% общего бюджета~--- закон Амдала: ускорять
то, на~что приходится 5~мс из~280, бессмысленно.

Сам риск-движок ускоряется без~криптографии: лимиты в~памяти вместо
запросов к~БД, кэш признаков клиента и~устройства, асинхронные проверки там,
где задержка ответа укладывается в~SLA. Это инженерия данных, не~криптография.

Карту по~протоколу ускорить тоже нельзя. Фаза~5~--- около 285~мс на~одну
\texttt{GENERATE~AC}~--- работа MPA (Mir Payment Application, апплет
на~чипе) на~ограниченном CPU карты.
Mir~Kernel, национальное L2-ядро НСПК на~терминале (в~EMVCo Book~C
  не~входит, работает параллельно международным Kernel~2~= M/Chip~Contactless
и~Kernel~3~= qVSDC поверх общего EMV-фреймворка), уже использует один
GENERATE~AC без~Issuer Script Processing. То, что Quick~Chip и~M/Chip~Fast
(см.~раздел~\ref{sec:emv-fast-flows}) делают для~контактного dip'a
на~уровне терминального ядра, в~бесконтактном пути встроено по~дизайну. Сократить
эти 285~мс можно только сменой поколения чипа с~аппаратным криптомодулем~---
то есть массовым перевыпуском карт.

В~сети между картой и~эмитентом остаётся одна честная экономия: прямое
соединение шлюза эквайера с~ОПКЦ, без~промежуточных интеграторов,~---
10--30~мс на~переходах.

Что не~ускоряется ничем: физика. Сигнал между ЦОДами в~разных регионах России
не~уложится меньше чем в~20--40~мс на~round trip, и~если эмитент~--- в~Москве,
а~эквайер~--- во~Владивостоке, авторизация физически не~будет быстрой.
Внутрироссийский маршрут спасает от~межконтинентальных задержек, но~не
от~собственной географии.

\section{Что осталось за~этой цепочкой}\label{sec:one-tap-out-of-scope}

В~описанной транзакции всё прошло штатно. Реальный поток платежей включает также
отказы, переключения \textit{fallback} и~частичные сбои, которые
в~эту главу не~вошли:

\emph{Decline эмитента.} Если фаза~9 завершается \texttt{Response~Code~05}
(do not honor), \texttt{14} (invalid card number) или другим decline-кодом~---
терминал получает короткое MTI~0110 и~издаёт другой beep (или показывает
сообщение). Логика обработки decline-кодов и~реакции эквайера~---
раздел~\ref{sec:iso8583-hex-dump}.

\emph{Сервис Резервной Авторизации ОПКЦ.} Если эмитент недоступен, ОПКЦ сам
принимает решение по~правилам Резервной Авторизации (раздел~13 Правил
ПС~Мир)~\cite{nspk-mir-rules-v43}. Карта получает approve или decline
без~участия эмитента; уведомление эмитенту уйдёт постфактум. Этот механизм
работает целиком внутри ОПКЦ и~не~пересекается с~Visa STIP или Mastercard stand-in.

\emph{Кобейджинг Мир--UnionPay.} Если карта~--- кобейджинговая Мир--UnionPay
и~операция проходит за~рубежом, транзакция маршрутизируется не~через ОПКЦ, а~через сеть
UnionPay International. На~2026~год эмиссия Мир--UnionPay ограничена несколькими банками,
а~зарубежный приём резко сужен после SDN-листинга НСПК. Подробно~---
раздел~\ref{sec:mir-history}.

\emph{Fallback на~контактный чип.} Если бесконтактный обмен не~завершился
(карта вышла из~поля, FWT превышен), терминал переключается на~контактный
интерфейс: карта вставляется в~слот. APDU те~же, физика другая
(ISO~7816 вместо ISO~14443), бюджет вырастает до~3--5~секунд~---
раздел~\ref{sec:emv-fallback}.

\emph{Token вместо PAN.} Если транзакция~--- через Mir~Pay, в~DE~2 уходит
не~PAN, а~токен НСПК. Сервис токенизации НСПК (раздел~\ref{sec:mir-pay})
делает де-токенизацию на~своей стороне и~передаёт эмитенту настоящий PAN;
ARQC при~этом считается ключом DPAN, а~не PAN-ключом. Apple~Pay и~Google~Pay
для~российских карт с~2022~года недоступны~\cite{nfcw-applepay-googlepay-mir-2022};
оплата телефоном в~России~--- это Mir~Pay на~Android с~HCE. Подробно~---
гл.~\ref{ch:tokenization}.

\emph{3-D Secure.} В~бесконтактной транзакции 3DS не~участвует~--- держатель
физически присутствует у~терминала, и~аутентификация уже выполнена через
CVM\,/\,CDCVM на~стороне карты или устройства. 3DS включается только в~CNP-сценарии;
для~карт \enquote{Мир} это MirAccept, гл.~\ref{ch:3ds} и~раздел~\ref{sec:mir-accept}.

\emph{Chargeback.} Если держатель оспорит операцию через несколько дней или
недель, начинается параллельная цепочка~--- по~правилам ПС~Мир,
доказательства, арбитраж НСПК. Эта цепочка идёт за~пределами 1200~мс
tap-to-beep и~разбирается в~главе~\ref{ch:disputes}.
